\documentclass[11pt]{article}
\usepackage{amsmath, amssymb}
\usepackage[margin=1in]{geometry}
\usepackage{fancyhdr}
\usepackage{hyperref}

\hypersetup{
    colorlinks=true,
    linkcolor=blue,
    urlcolor=blue,
    citecolor=green
}

% ---- Fill these in ----
\newcommand{\courseName}{NEU 314}
\newcommand{\taName}{Keith Murray}
\newcommand{\sectionInfo}{Precept B01}
% ------------------------

\title{Precept: Bias and Variance of an Estimator}
\date{\today}

\pagestyle{fancy}
\fancyhf{}
\lhead{\courseName}
\chead{\sectionInfo}
\rhead{\today}
\cfoot{\thepage}

\begin{document}

\maketitle
\thispagestyle{fancy}

\section*{Problem}
Recall the two properties of an estimator $\hat{\theta}$ from lecture, where the expectation $\mathbb{E}$ is an average over draws of the dataset $m$:
\[
\textbf{Bias: } b(\theta)=\mathbb{E}[\hat{\theta}]-\theta
\qquad\qquad
\textbf{Variance: } \mathrm{var}(\theta)=\mathbb{E}\Big[\big(\hat{\theta}-\mathbb{E}[\hat{\theta}]\big)^2\Big]
\]
In lecture we saw that the estimator $\hat{\theta}(m)=7$, which ignores the data entirely, has variance $0$ and bias $7-\theta$. Now let's look at an estimator that actually uses the data.

\begin{enumerate}
    \item \textbf{The single-trial estimator.} You present the same stimulus to a neuron on $n$ trials and record the spike counts $m=\{r_1, r_2, \ldots, r_n\}$. The trials are independent, and each spike count has mean $\theta$ and variance $\sigma^2$:
    \[
    \mathbb{E}[r_i]=\theta, \qquad \mathrm{Var}(r_i)=\sigma^2 \qquad \text{for every } i
    \]
    A lazy experimenter decides to estimate $\theta$ using only the first trial:
    \[
    \hat{\theta}(m)=r_1
    \]
    \begin{enumerate}
        \item What is the bias of this estimator? Is it unbiased?
        \item What is the variance of this estimator?
        \item The experimenter records $n=1000$ trials instead of $n=10$. How do the bias and variance of $\hat{\theta}=r_1$ change? What does this tell you about the estimator?
        \item Compare $\hat{\theta}=r_1$ to the constant estimator $\hat{\theta}=7$. Each one gets one property ``right.'' Which property does each get right, and what does each give up?
    \end{enumerate}
\end{enumerate}

\newpage
\section*{Solution}

\subsection*{Part (a)}
The expected value of the estimator is just the expected value of the first spike count:
\[
\mathbb{E}[\hat{\theta}]=\mathbb{E}[r_1]=\theta
\]
so the bias is
\[
b(\theta)=\mathbb{E}[\hat{\theta}]-\theta=\theta-\theta=0
\]
The estimator is unbiased, for \emph{every} value of $\theta$. On average over many repeated experiments, the first trial lands right on the true mean.

\subsection*{Part (b)}
Using $\mathbb{E}[\hat{\theta}]=\theta$ from Part (a),
\[
\mathrm{var}(\theta)=\mathbb{E}\Big[\big(\hat{\theta}-\mathbb{E}[\hat{\theta}]\big)^2\Big]=\mathbb{E}\big[(r_1-\theta)^2\big]=\sigma^2
\]
The last step is just the definition of the variance of $r_1$. The estimator is exactly as noisy as a single spike count.

\subsection*{Part (c)}
Nothing changes! The bias is still $0$ and the variance is still $\sigma^2$, because the estimator never looks at $r_2, \ldots, r_n$. The experimenter paid for 999 extra trials and threw all of them away. Being unbiased does not make an estimator good: a better estimator should get more precise as we collect more data. (For example, the sample mean $\bar{r}=\frac{1}{n}\sum_i r_i$ is also unbiased, but its variance is $\sigma^2/n$, which shrinks as $n$ grows.)

\subsection*{Part (d)}
\begin{center}
\begin{tabular}{lcc}
 & Bias & Variance \\
\hline
$\hat{\theta}=7$ & $7-\theta$ & $0$ \\
$\hat{\theta}=r_1$ & $0$ & $\sigma^2$
\end{tabular}
\end{center}
The constant estimator gets the variance right: it gives the same answer every time, so it has no variance at all. But it ignores the data, so it is biased unless the true value happens to be exactly $7$. The single-trial estimator gets the bias right: it is centered on the truth. But it inherits all of the trial-to-trial noise of one spike count. Bias and variance are separate properties, and an estimator can do well on one while doing poorly on the other.

\end{document}